\section*{第 \ref{ch:g7:atoms-and-molecules} 章——原子与分子}

\begin{solution}{exo:g7:atoms-and-molecules:1}
\omterm{def:g7:atoms-and-molecules:atom}{原子}是构成一切物质的极微小粒子之一，
（已知约有一百种）。\omterm{def:g7:atoms-and-molecules:molecule}{分子}是由若干\omterm{def:g7:atoms-and-molecules:atom}{原子}
紧密结合在一起的集团，是保持该物质
性质的最小粒子。
\end{solution}

\begin{solution}{exo:g7:atoms-and-molecules:2}
\ce{H}, \ce{C}, \ce{N}, \ce{O}, \ce{Na}, \ce{Fe}.
\end{solution}

\begin{solution}{exo:g7:atoms-and-molecules:3}
1 个碳\omterm{def:g7:atoms-and-molecules:atom}{原子}和 2 个氧\omterm{def:g7:atoms-and-molecules:atom}{原子}，共 3 个\omterm{def:g7:atoms-and-molecules:atom}{原子}。
\end{solution}

\begin{solution}{exo:g7:atoms-and-molecules:4}
水、氧气、甲烷、氮气。
\end{solution}

\begin{solution}{exo:g7:atoms-and-molecules:5}
\ce{NO2}.
\end{solution}

\begin{solution}{exo:g7:atoms-and-molecules:6}
\ce{O2} 是一个\omterm{def:g7:atoms-and-molecules:molecule}{分子}，由两个结合在一起的氧\omterm{def:g7:atoms-and-molecules:atom}{原子}构成；\ce{2O} 是两个
彼此分开的氧\omterm{def:g7:atoms-and-molecules:atom}{原子}。\ce{CO} 是由一个碳\omterm{def:g7:atoms-and-molecules:atom}{原子}和一个氧\omterm{def:g7:atoms-and-molecules:atom}{原子}组成的\omterm{def:g7:atoms-and-molecules:molecule}{分子}
（两个大写字母，两个\omterm{def:g7:atoms-and-molecules:symbol}{符号}）；\ce{Co} 是一个钴\omterm{def:g7:atoms-and-molecules:atom}{原子}
（一个\omterm{def:g7:atoms-and-molecules:symbol}{符号}，大写字母后接小写字母）。
\end{solution}

\begin{solution}{exo:g7:atoms-and-molecules:7}
\ce{3H2O}：$3\times 2 = 6$ 个氢\omterm{def:g7:atoms-and-molecules:atom}{原子}，$3\times 1 = 3$ 个氧
\omterm{def:g7:atoms-and-molecules:atom}{原子}。\ce{5O2}：不含氢，$5\times 2 = 10$ 个氧\omterm{def:g7:atoms-and-molecules:atom}{原子}。
\end{solution}

\begin{solution}{exo:g7:atoms-and-molecules:8}
黑球是碳，红球是氧：一个碳、两个氧，
即 \ce{CO2}，二氧化碳。
\end{solution}

\begin{solution}{exo:g7:atoms-and-molecules:9}
道尔顿所认为的水是 \ce{HO}。真正的\omterm{def:g7:atoms-and-molecules:molecule}{分子} \ce{H2O} 比它多
一个氢\omterm{def:g7:atoms-and-molecules:atom}{原子}：道尔顿的\omterm{def:g7:atoms-and-molecules:molecule}{分子}缺少了一个\omterm{def:g7:atoms-and-molecules:atom}{原子}。
\end{solution}

\begin{solution}{exo:g7:atoms-and-molecules:10}
纯水不是\omterm{def:g6:pure-substances-and-mixtures:mixture}{混合物}：它的\omterm{def:g7:atoms-and-molecules:molecule}{分子}全都是相同的 \ce{H2O}
\omterm{def:g7:atoms-and-molecules:molecule}{分子}。\omterm{def:g4:air-a-mixture-of-gases:air}{空气}是\omterm{def:g6:pure-substances-and-mixtures:mixture}{混合物}：其中含有氮气\omterm{def:g7:atoms-and-molecules:molecule}{分子}、氧气
\omterm{def:g7:atoms-and-molecules:molecule}{分子}、氩\omterm{def:g7:atoms-and-molecules:atom}{原子}和少量二氧化碳\omterm{def:g7:atoms-and-molecules:molecule}{分子}。
\end{solution}

\begin{solution}{exo:g7:atoms-and-molecules:11}
一千万个\omterm{def:g7:atoms-and-molecules:atom}{原子}排成 \qty{1}{mm}，所以 \qty{1}{cm} $= \qty{10}{mm}$
需要 $10 \times$ 一千万 $=$ 一亿个\omterm{def:g7:atoms-and-molecules:atom}{原子}。
$\qty{100}{m} = \num{100000}$~mm 需要 $\num{100000}\times$ 一千万
$=$ 一万亿个\omterm{def:g7:atoms-and-molecules:atom}{原子}。
\end{solution}

\begin{solution}{exo:g7:atoms-and-molecules:12}
$6 + 12 + 6 = 24$ 个\omterm{def:g7:atoms-and-molecules:atom}{原子}。氢：$\frac{12}{24} = \qty{50}{\percent}$。
碳：$\frac{6}{24} = \qty{25}{\percent}$（氧占另外的
\qty{25}{\percent}）。
\end{solution}

\begin{solution}{pb:g7:atoms-and-molecules:1}
\textbf{1.} 氮气 \ce{N2}，2 个\omterm{def:g7:atoms-and-molecules:atom}{原子}；氧气 \ce{O2}，2 个\omterm{def:g7:atoms-and-molecules:atom}{原子}；二氧化碳
\ce{CO2}，3 个\omterm{def:g7:atoms-and-molecules:atom}{原子}。

\textbf{2.} 氩\omterm{def:g7:atoms-and-molecules:atom}{原子}不会结合成\omterm{def:g7:atoms-and-molecules:molecule}{分子}：氩由
单个\omterm{def:g7:atoms-and-molecules:atom}{原子}构成。

\textbf{3.} 氮：$78 \times 2 = 156$ 个\omterm{def:g7:atoms-and-molecules:atom}{原子}。氧：$21\times 2 =
42$ 个\omterm{def:g7:atoms-and-molecules:atom}{原子}。

\textbf{4.} $156 + 42 + 1 = 199$ 个\omterm{def:g7:atoms-and-molecules:atom}{原子}。

\textbf{5.} $\frac{156}{199} \approx 0.78$，约为 \qty{78}{\percent}。

\textbf{6.} 氧\omterm{def:g7:atoms-and-molecules:molecule}{分子}中：$16\times 2 = 32$；二氧化碳\omterm{def:g7:atoms-and-molecules:molecule}{分子}中：
$4 \times 2 = 8$；合计 $32 + 8 = 40$。

\textbf{7.} $42 - 40 = 2$ 个氧\omterm{def:g7:atoms-and-molecules:atom}{原子}不见了，相当于
一个氧\omterm{def:g7:atoms-and-molecules:molecule}{分子}的量。它们进入了水\omterm{def:g7:atoms-and-molecules:molecule}{分子} \ce{H2O}，人体用
食物中的氢制造这种水，其中一部分以水蒸气形式随呼出的气体排出
（图中已略去）。

\textbf{8.} 吸入：0。呼出：4（每个 \ce{CO2} 中各一个）。

\textbf{9.} 人体所用的食物中，糖类和脂肪含有碳
\omterm{def:g7:atoms-and-molecules:atom}{原子}，人体以呼出的二氧化碳形式将其释放。

\textbf{10.} $21 - 16 = 5$ 个氧\omterm{def:g7:atoms-and-molecules:molecule}{分子}消失了，并
出现了 4 个二氧化碳\omterm{def:g7:atoms-and-molecules:molecule}{分子}。

\textbf{11.} 氮气：两个蓝球由三根短棍相连。氧气：两个
红球由两根短棍相连。二氧化碳：一个黑球位于两个
红球之间，每个红球各以两根短棍与它相连，排成一条直线。氩：一个
浅蓝色的单球。

\textbf{12.} $156 + 32 + 12 + 1 = 201$ 个\omterm{def:g7:atoms-and-molecules:atom}{原子}（4 个二氧化碳\omterm{def:g7:atoms-and-molecules:molecule}{分子}
共含 12 个\omterm{def:g7:atoms-and-molecules:atom}{原子}）。这比吸入的\omterm{def:g4:air-a-mixture-of-gases:air}{空气}多 $201 - 199 = 2$ 个
\omterm{def:g7:atoms-and-molecules:atom}{原子}：从食物中获得 4 个碳\omterm{def:g7:atoms-and-molecules:atom}{原子}，减去随水
离开的 2 个氧\omterm{def:g7:atoms-and-molecules:atom}{原子}。
\end{solution}
